% Generated draft: requires mathematical and source review.
\documentclass[11pt]{article}
\usepackage[margin=25mm]{geometry}
\usepackage{fontspec}
\setmainfont{Latin Modern Roman}
\usepackage{amsmath,amssymb,mathtools,unicode-math}
\setmathfont{Latin Modern Math}
\usepackage{enumitem,needspace,xurl,hyperref}
\setlength{\parindent}{0pt}
\setlength{\parskip}{5pt}
\setlength{\emergencystretch}{3em}
\newcommand{\N}{\mathbb{N}}
\newcommand{\Z}{\mathbb{Z}}
\newcommand{\Q}{\mathbb{Q}}
\newcommand{\R}{\mathbb{R}}
\newcommand{\C}{\mathbb{C}}
\newcommand{\F}{\mathbb{F}}
\newcommand{\A}{\mathbb{A}}
\newcommand{\PP}{\mathbb{P}}
\newcommand{\E}{\mathbb{E}}
\newcommand{\eps}{\varepsilon}
\newcommand{\dd}{\,\mathrm{d}}
\begin{document}
\section*{UnsolvedMath ID: 3700020}
\textbf{Source problem code:} catalogue entry
\subsection{Definitions and mathematical statement}
Identify the plane with \(\C\), and write \(D_r=\{z\in\C:|z|\le r\}\) and \(C_r=\{z\in\C:|z|=r\}\) for \(r>0\). The question concerns a positive Borel measure \(\mu\) satisfying the growth condition
\[
\exists c>0\ \exists\alpha\in(0,1/2)\ \forall r>0,\qquad \mu(D_r)\le cr^\alpha,
\]
and its logarithmic potential
\[
u(z)=\int_{\C}\log\left|1-\frac z\zeta\right|\,\mathrm d\mu(\zeta).
\]
The all-radii growth interpretation implies \(\mu(\{0\})=0\), so the integrand at \(\zeta=0\) may be ignored. Local constancy is understood for finite real values of the potential, rather than a constant value \(-\infty\).

Does there exist such a measure and an open set \(U\subseteq\C\) such that
\[
\forall r>0,\quad U\cap C_r\ne\varnothing,
\]
\[
\forall z\in U\ \exists\varepsilon>0,\quad
B(z,\varepsilon)\subseteq U\quad\text{and}\quad
\forall w\in B(z,\varepsilon),\ u(w)=u(z)\in\R?
\]
Local constancy permits different constants on different components of \(U\); no connectedness or single common value is imposed. The source does not explicitly exclude \(\mu=0\). If positive means nonnegative, that measure gives a trivial affirmative example. The substantive nonzero-measure interpretation adds \(\mu(\C)>0\), and this missing qualification must be checked rather than silently inserted. The source also leaves the radius quantifier implicit; the displayed formulation uses all \(r>0\), not merely sufficiently large radii. The supplied research summary reports no located resolution and attributes the problem to Eremenko--Lewis; this is an unverified literature assessment.
\subsection{Short English statement}
Can a nontrivial positive plane measure of growth \(O(r^\alpha)\), \(0<\alpha<1/2\), have a logarithmic potential locally constant on an open set intersecting every centered circle? Nontriviality and the radius convention need source confirmation.
\subsection{Sources}
\begin{itemize}
\item Local UnsolvedMath registry, permanent ID \(3700020\), exact source code \texttt{catalogue entry}, ``Locally constant logarithmic potentials.'' Registry proposer: Alexandre Eremenko and cited proposers; proposed year 2008.
\item Eremenko, \emph{My favorite unsolved problems}, Potential theory, Problem 1; \url{https://www.math.purdue.edu/~eremenko/uns1.html}. Recorded access 2026-07-29; extraction recorded as pdf-text.
\item Supplied OPEN-TRIAGE research summary attributes the problem to Eremenko--Lewis and names \texttt{equivpot.pdf}, but supplies no complete publication citation. Its explanation of the exponent threshold is not accompanied by a theorem or reference.
\end{itemize}
\end{document}
